\section{Implementation}
\label{sec:impl}

Pigs is implemented as a Rust workspace with 14 crates, organized in five layers:
Pigs 以 Rust workspace 形式实现，包含 14 个 crate，组织为五个层次：

\begin{center}
\begin{tabular}{lll}
\toprule
Layer & Crates & Role \\
\midrule
0 & pigs-core & Core types: Message, ContentBlock, ApiClient trait \\
1 & pigs-config, pigs-permissions, & Config, permissions, session, prompts \\
  & pigs-session, pigs-prompts & \\
2 & pigs-llm, pigs-tools, pigs-mcp & LLM clients, built-in tools, MCP client \\
3 & pigs-api, pigs-proxy & Phased runtime, HTTP proxy + loopback \\
4 & pigs-cli, pigs-tui, pigs & Agent logic, terminal UI, product binary \\
\bottomrule
\end{tabular}
\end{center}

\subsection{Protocol Codec}

The \texttt{ProtocolCodec} in \texttt{pigs-api} handles parsing, mutation, and response extraction for all three protocols. It validates the ``current user input'' (the last user message, not just the last non-system message) and identifies trailing tool-result groups for continuation matching.
\texttt{pigs-api} 中的 \texttt{ProtocolCodec} 负责所有三种协议的解析、修改和响应提取。它验证"当前用户输入"（最后一条 user 消息，而非仅仅是最后一条非系统消息），并识别尾部的工具结果组以进行续传匹配。

Each protocol has dedicated parsing logic:
每种协议具有专用的解析逻辑：
\begin{itemize}[nosep]
    \item \textbf{Chat}: messages array with role-based content; tool calls as \texttt{tool\_calls} array; tool results as \texttt{role: "tool"} messages.
    \textbf{Chat}：具有基于角色内容的 messages 数组；工具调用为 \texttt{tool\_calls} 数组；工具结果为 \texttt{role: "tool"} 消息。
    \item \textbf{Anthropic}: messages with content block arrays; tool calls as \texttt{tool\_use} blocks; tool results as \texttt{tool\_result} blocks within user messages.
    \textbf{Anthropic}：具有内容块数组的 messages；工具调用为 \texttt{tool\_use} 块；工具结果为 user 消息内的 \texttt{tool\_result} 块。
    \item \textbf{Responses}: input array with typed items; tool calls as \texttt{function\_call} items; tool results as \texttt{function\_call\_output} items.
    \textbf{Responses}：具有类型化条目的 input 数组；工具调用为 \texttt{function\_call} 条目；工具结果为 \texttt{function\_call\_output} 条目。
\end{itemize}

The codec clones the original body and modifies only the JSON path leading to the current user input, leaving all other fields structurally identical.
编解码器克隆原始 body 并仅修改指向当前用户输入的 JSON 路径，保持所有其他字段在结构上完全一致。

\subsection{HTTP Runtime}

The \texttt{HttpPhasedRuntime} drives the phase loop. For each phase, it:
\texttt{HttpPhasedRuntime} 驱动阶段循环。对于每个阶段，它：
\begin{enumerate}[nosep]
    \item Constructs a phase request via \texttt{phase\_request()} (Pre/Executor/Post prompt + transcript).
    通过 \texttt{phase\_request()} 构建阶段请求（Pre/Executor/Post 提示 + 记录）。
    \item Sends the request via \texttt{PhaseTransport::send\_streaming()} (streaming) or \texttt{send()} (non-streaming).
    通过 \texttt{PhaseTransport::send\_streaming()}（流式）或 \texttt{send()}（非流式）发送请求。
    \item Extracts the normalized model output (visible text, tool calls, usage, stop reason).
    提取归一化的模型输出（可见文本、工具调用、使用量、停止原因）。
    \item Checks for tool calls $\rightarrow$ pauses and stores continuation.
    检查工具调用 $\rightarrow$ 暂停并存储续传。
    \item Advances the state machine via \texttt{OrchestrationState::advance()}.
    通过 \texttt{OrchestrationState::advance()} 推进状态机。
\end{enumerate}

The runtime is transport-agnostic: it depends only on the \texttt{PhaseTransport} trait, not on HTTP or axum. This allows the same runtime to be used with any transport implementation.
运行时与传输无关：它仅依赖于 \texttt{PhaseTransport} trait，不依赖 HTTP 或 axum。这使得同一运行时可与任何传输实现配合使用。

\subsection{Loopback Transport}

The \texttt{LoopbackPhaseTransport} in \texttt{pigs-proxy} sends phase subrequests via HTTP to the local proxy. It parses both non-streaming JSON and SSE responses from the upstream provider, aggregating:
\texttt{pigs-proxy} 中的 \texttt{LoopbackPhaseTransport} 通过 HTTP 将阶段子请求发送至本地代理。它解析来自上游提供商的非流式 JSON 和 SSE 响应，聚合：
\begin{itemize}[nosep]
    \item Chat: text deltas, tool call deltas (merged by index), finish reason, usage, provider extensions.
    Chat：文本增量、工具调用增量（按索引合并）、完成原因、使用量、提供商扩展。
    \item Anthropic: content blocks (text, thinking, signature, tool\_use with \texttt{input\_json\_delta} accumulation), stop reason, usage.
    Anthropic：内容块（text、thinking、signature、tool\_use 及 \texttt{input\_json\_delta} 累积）、停止原因、使用量。
    \item Responses: output items (message, reasoning, function\_call), status, incomplete details, usage.
    Responses：输出条目（message、reasoning、function\_call）、状态、不完整详情、使用量。
\end{itemize}

Malformed streams (missing termination events, unknown content block indices, open blocks at stream end) are reported as typed errors, not silently accepted.
格式错误的流（缺少终止事件、未知内容块索引、流结束时未关闭的块）被报告为类型化错误，而非被静默接受。

\subsection{Body Cleaning}

The proxy's \texttt{clean\_empty\_messages} function removes messages with empty content, but preserves assistant messages that contain \texttt{tool\_calls} (even if their content is empty). For messages with empty content and tool calls, the content is replaced with a single space to satisfy upstream providers that reject empty strings.
代理的 \texttt{clean\_empty\_messages} 函数移除内容为空的消息，但保留包含 \texttt{tool\_calls} 的 assistant 消息（即使其内容为空）。对于内容为空且含工具调用的消息，内容被替换为单个空格，以满足上游提供商对空字符串的拒绝要求。

\subsection{Agent Client and Terminal UI}

Pigs ships with a general-purpose agent client (similar to PI~\citep{omo2026} or OpenCode), which connects to the local API server via HTTP. The key design point is that the client can use either model identifier: \texttt{\{model\}} (passthrough, no orchestration) or \texttt{\{model\}-pig} (phased orchestration). This allows the same agent to toggle orchestration on or off simply by changing the model name in the API request, with no other code changes.
Pigs 附带一个通用 Agent 客户端（类似 PI~\citep{omo2026} 或 OpenCode），通过 HTTP 连接本地 API 服务器。关键设计在于：客户端可以使用任一模型标识符——\texttt{\{model\}}（透传，无编排）或 \texttt{\{model\}-pig}（相位编排）。这使得同一 Agent 仅通过更改 API 请求中的模型名即可开启或关闭编排，无需其他代码变更。

The user interface is a full-screen terminal UI built with \textbf{ratatui} and \textbf{crossterm}, implemented in the \texttt{pigs-tui} crate. The layout follows PI's convention, with five vertically-stacked regions: (1) a header bar showing the version, model, and keybinding hints; (2) a scrollable chat history rendering user messages, assistant responses (with Markdown formatting via \texttt{pulldown-cmark}), tool call results, and bash output; (3) a status line with a working spinner; (4) a multi-line text editor (Emacs-style keybindings via \texttt{tui-textarea}); and (5) a footer status bar displaying the working directory, git branch, token statistics, and current model.
用户界面为全屏终端 UI，使用 \textbf{ratatui} 和 \textbf{crossterm} 实现，位于 \texttt{pigs-tui} crate 中。布局遵循 PI 的约定，包含五个垂直堆叠区域：(1) 顶部栏显示版本、模型和快捷键提示；(2) 可滚动的聊天历史区，渲染用户消息、助手回复（通过 \texttt{pulldown-cmark} 进行 Markdown 格式化）、工具调用结果和 bash 输出；(3) 状态行显示工作旋转指示器；(4) 多行文本编辑器（通过 \texttt{tui-textarea} 实现 Emacs 风格快捷键）；(5) 底部状态栏显示工作目录、git 分支、token 统计和当前模型。

Streaming output is handled through an event broker architecture: the agent's \texttt{run\_turn\_with\_callback} method forwards \texttt{TextDelta}, \texttt{ToolStart}, and \texttt{ToolEnd} events from the phased runtime's \texttt{ProgressSink} to the TUI via a \texttt{tokio::mpsc} channel. This enables real-time, token-by-token rendering of the model's response, live tool execution display, and phase-aware status updates. The TUI gracefully falls back to a line-based REPL (via \texttt{rustyline}) if terminal initialization fails.
流式输出通过事件代理架构处理：Agent 的 \texttt{run\_turn\_with\_callback} 方法将相位运行时 \texttt{ProgressSink} 的 \texttt{TextDelta}、\texttt{ToolStart} 和 \texttt{ToolEnd} 事件通过 \texttt{tokio::mpsc} 通道转发到 TUI。这实现了模型响应的实时逐 token 渲染、工具执行的实时显示和相位感知状态更新。如果终端初始化失败，TUI 会优雅地回退到基于行的 REPL（通过 \texttt{rustyline}）。

The built-in tool set is aligned with PI: \texttt{read}, \texttt{bash}, \texttt{edit}, and \texttt{write} are enabled by default, with optional \texttt{grep}, \texttt{find}, and \texttt{ls}. Configuration paths follow PI's convention: \texttt{\textasciitilde/.pig/} for global state (sessions, logs, skills, config) and \texttt{.pig/} for project-local overrides.
内置工具集与 PI 对齐：\texttt{read}、\texttt{bash}、\texttt{edit} 和 \texttt{write} 默认启用，可选启用 \texttt{grep}、\texttt{find} 和 \texttt{ls}。配置路径遵循 PI 的约定：\texttt{\textasciitilde/.pig/} 用于全局状态（会话、日志、技能、配置），\texttt{.pig/} 用于项目级覆盖。
